% =====================================================================
%  Originality of the results --- an index
%  Standalone: latexmk -pdf standalone.tex
%  In-thesis:  \subimport{originality_index/}{originality.tex}
% =====================================================================

% --- Guarded setup: harmless if the thesis preamble already defines these ---
\providecommand{\yp}{\textit{Yersinia pestis}}
\providecommand{\ft}{\textit{Francisella tularensis}}

\makeatletter
\@ifundefined{oi@colours}{%
  \definecolor{oiink}{HTML}{1A1C1F}
  \definecolor{oisoft}{HTML}{565B62}
  \definecolor{oiaccent}{HTML}{0072B2}
  \definecolor{oipale}{HTML}{EFF6FB}
  \definecolor{oiedge}{HTML}{BBD6E9}
  \definecolor{oinewfill}{HTML}{E2EFF9}
  \definecolor{oigrayfill}{HTML}{F0F0F0}
  \definecolor{oigrayedge}{HTML}{C8C8C8}
  \definecolor{oirule}{HTML}{D8DDE2}
  \gdef\oi@colours{done}%
}{}
\makeatother

\providecommand{\Y}{}
\newcolumntype{Z}{>{\raggedright\arraybackslash}X}
\newcolumntype{W}{>{\raggedright\arraybackslash}p{20mm}}
\newcolumntype{N}{>{\raggedright\arraybackslash}p{36mm}}
\newcommand{\oihead}[1]{\textcolor{oisoft}{\footnotesize\bfseries\MakeUppercase{#1}}}
\newcommand{\anchor}[1]{\textcolor{oiaccent}{\textbf{#1}}}
\newcommand{\oisec}[1]{%
  \par\addvspace{0.7\baselineskip}%
  \noindent{\oihead{#1}}\par
  \nobreak\vspace{-0.15\baselineskip}%
  {\color{oirule}\rule{\linewidth}{0.5pt}}\par
  \nobreak\vspace{0.1\baselineskip}}

% =====================================================================
\clearpage
\phantomsection
\addcontentsline{toc}{chapter}{Originality of the results}
\markboth{ORIGINALITY OF THE RESULTS}{ORIGINALITY OF THE RESULTS}

{\LARGE\bfseries Originality of the results\par}
\vspace{0.35\baselineskip}
{\color{oisoft}An index of what is inherited, what is extended, and what is new\par}
\vspace{0.3\baselineskip}
{\color{oirule}\rule{\linewidth}{0.8pt}}
\vspace{0.1\baselineskip}

\noindent
Every entry names the published ancestor and the increment made to it, and points to the
theorem and page where the result is stated. The analytical spine of this thesis --- the
linear birth--death process with per-capita total catastrophe --- is not new, and is not
claimed to be: it is Karlin and Tavar\'e's killed birth--death process, and the constant
$A(p)$ of \hyperref[gw:ch:galton-watson]{Chapter~3} is Kolmogorov's. Built on that spine
are a closed-form intracellular calculus, its embedding in a within-host renewal system,
a two-type deformation, and a criterion that orders bursting against budding. Four
statements were not found in the literature as written; they come first, each beside the
neighbour that comes closest.

% ---------------------------------------------------------------------
\vspace{0.2\baselineskip}
\begin{center}
\begin{tikzpicture}[
  band/.style   = {draw, rounded corners=2pt, line width=0.5pt,
                   text width=115mm, inner xsep=3.4mm, inner ysep=1.9mm,
                   align=left, font=\footnotesize},
  inherit/.style= {band, fill=oigrayfill, draw=oigrayedge, text=oisoft},
  built/.style  = {band, fill=oipale, draw=oiedge},
  newband/.style= {band, fill=oinewfill, draw=oiaccent, line width=0.9pt},
  tag/.style    = {font=\scriptsize\bfseries, text=oisoft, anchor=east},
  node distance = 2.1mm,
]

\node[newband] (top) {%
  \textbf{\textcolor{oiaccent}{A criterion for release strategy.}}
  One dimensionless group $L$ orders bursting against budding at matched yield: the
  extinction gap is $(L-1)(1/m-1)$, the two offspring laws coincide at $L=1$, and the
  order reverses under a finite immune budget.};
\node[tag] at (top.west) [xshift=-1.8mm] {Ch 6, 8};

\node[built, below=of top] (mid) {%
  \textbf{Renewal embedding and a two-type deformation.}
  Closed-form burst kernels drive an age-structured within-host system; the triangular
  two-type process is solved under type-specific killing.};
\node[tag] at (mid.west) [xshift=-1.8mm] {Ch 6, 7};

\node[built, below=of mid] (calc) {%
  \textbf{An intracellular calculus.}
  The released count $W_t$ carried alongside the load $X_t$; every load moment organised
  as a polynomial in $I(t)$; burst-size and burst-time laws.};
\node[tag] at (calc.west) [xshift=-1.8mm] {Ch 4, 5};

\node[inherit, below=of calc] (base) {%
  \textbf{Inherited framework.}
  Linear birth--death with killing: linear-fractional generating function, geometric
  occupancy, geometric killing position (Karlin and Tavar\'e 1982). The constant $A(p)$
  (Kolmogorov 1938). Burst versus continuous release at matched $R_0$ (Pearson et al.\ 2011).};
\node[tag] at (base.west) [xshift=-1.8mm] {Ch 2--5};

\draw[-{Stealth[length=2.2mm]}, oiaccent, line width=0.7pt]
  ($(base.south west)+(-24mm,1mm)$) -- ($(top.north west)+(-24mm,-1mm)$);
\node[rotate=90, font=\scriptsize, text=oisoft, anchor=south]
  at ($(base.south west)!0.5!(top.north west)+(-25mm,0)$) {increasing ownership};

\end{tikzpicture}
\end{center}

% =====================================================================
\oisec{1. Statements not found in the literature as written}

{\footnotesize\noindent
The method behind each is published; these are original \emph{statements} obtained by
specialising it, not a new theory of branching. Each was searched against the chapter
bibliographies and the wider literature, then challenged adversarially.\par}

\vspace{0.35\baselineskip}
{\footnotesize
\begin{xltabular}{\linewidth}{@{}WZN@{}}
\toprule
\oihead{Where} & \oihead{The statement} & \oihead{Nearest published work} \\
\midrule
\endfirsthead
\toprule
\oihead{Where} & \oihead{The statement} & \oihead{Nearest published work} \\
\midrule
\endhead
\bottomrule
\endlastfoot

\anchor{Thm 6.8}\newline p.\,163 &
\textbf{The flooding identity.} At matched lifetime yield the extinction probabilities of
bursting and budding differ by exactly $\Delta z_{\mathrm{ext}}=(L-1)(1/m-1)$, so bursting
favours establishment precisely when $L>1$, equivalently when $1/\delta<1/\lambda+1/\mu$. &
Pearson, Krapivsky and Perelson (2011) find $p_v^{\mathrm{burst}}\le p_v^{\mathrm{cont}}$ ---
bursting always favoured --- from a fixed burst size, with no parameter that can reverse the
order. \\
\addlinespace[3pt]

\anchor{Prop 6.9}\newline p.\,165 &
\textbf{Coincidence of the offspring laws.} At $L=1$ the bursting and budding offspring
\emph{distributions} coincide, not merely their extinction probabilities: the burst
generating function collapses to the geometric $(1-\varrho)/(1-\varrho z)$. The variance
gap $2q^{2}V_\infty(L-1)/(a-1)$ changes sign with $L-1$. &
Pearson's geometric comparator; Karlin and Tavar\'e's geometric killing position. Neither
yields a coincidence of laws. \\
\addlinespace[3pt]

\anchor{Thm 7.4}\newline p.\,189 &
\textbf{Exact finite-time survival under two-type killing.} A closed form for the
no-catastrophe probability of the triangular two-type birth--death process with
type-specific catastrophe rate $\delta_1X+\delta_2Y$, as a Wronskian mixture of a
real-axis ${}_2F_1$ pair. &
Antal and Krapivsky (2011) apply the same Gauss reduction to terminal counts without
killing. Setting $\delta=0$ here forces $S\equiv G\equiv1$ and does \emph{not} return
their generating functions. \\
\addlinespace[3pt]

\anchor{Eq.\ 7.50}\newline p.\,201 &
\textbf{The killed bivariate generating function.} The defective PGF under the same
killing, by the same Wronskian mixture. This is the object for which $\delta_i=0$
\emph{does} recover Antal--Krapivsky, and it is what makes the previous entry a different
observable rather than a reparametrisation. &
Antal and Krapivsky (2011), eqs.\ (25)--(29), recovered here as the $\delta_i\to0$ limit. \\
\end{xltabular}\par}

% =====================================================================
\oisec{2. Extensions: an ancestor, and what is added to it}

{\footnotesize\noindent
Ordered by chapter. In each case a published result is the ancestor; the thesis changes
the process, the object, the hypotheses, or the biological reading.\par}

\vspace{0.35\baselineskip}
{\footnotesize
\begin{xltabular}{\linewidth}{@{}WZN@{}}
\toprule
\oihead{Where} & \oihead{The increment} & \oihead{Ancestor} \\
\midrule
\endfirsthead
\toprule
\oihead{Where} & \oihead{The increment} & \oihead{Ancestor} \\
\midrule
\endhead
\bottomrule
\endlastfoot

\anchor{\S2.2.4.1}\newline p.\,35 &
Gated release in which the opening intensity $\eta A_t$ depends on the accumulated load,
rather than switching autonomously. &
Davis's piecewise-deterministic Markov processes; autonomous on--off queues. \\
\addlinespace[3pt]

\anchor{Thm 3.5}\newline p.\,82 &
$A(p)=2\psi_{2p}(1/2)$ identifies Kolmogorov's constant with the Koenigs linearising
function of the logistic map, and that function is hypertranscendental --- so it admits no
elementary closed form. &
Kolmogorov (1938) for $A(p)$ itself; Koenigs (1884); Becker and Bergweiler (1995) for
hypertranscendence. \\
\addlinespace[3pt]

\anchor{Thms 4.13--4.14}\newline pp.\,100--101 &
The released count $W_t$ carried as a second coordinate: its mean, second moment and
variance at \emph{finite} time, as functions of $I(t)$. &
Karlin and Tavar\'e's killing position is the \emph{eventual} geometric, not a
time-dependent release coordinate. \\
\addlinespace[3pt]

\anchor{Prop 4.17}\newline p.\,103 &
Every load moment is a polynomial $Q_p(I)$ of degree $p+1$ with roots $a,b$ and leading
coefficient $(-1)^p p!\,(\lambda/\delta)^p$ --- one economy organising the whole moment
hierarchy. &
The moments themselves sit inside the 1982 generating function; the organisation is what
is added. \\
\addlinespace[3pt]

\anchor{Thm 5.10}\newline p.\,126 &
Burst size equals the quasi-stationary law, proved by size-biasing: the weight $k$ makes
the burst-time integrand a perfect differential, which telescopes to the endpoint.
Remark 5.11 gives a \emph{sufficient} criterion for it; necessity is left open. &
Karlin and Tavar\'e give the killing-position law and the quasi-stationary law separately,
in different lemmas. \\
\addlinespace[3pt]

\anchor{Props 5.12--5.13}\newline pp.\,130--131 &
Burst timing: at $\mu=0$ the time to a burst of size $n$ is harmonic in $n$ with a Gumbel
limit; with death, a closed mean conditional on bursting. &
Classical Yule maxima, not previously written as a law for a lytic single-cell assay. \\
\addlinespace[3pt]

\anchor{\S5.9--5.10}\newline pp.\,134--147 &
Multi-founder kernels $g_k$ with sub-extensive yield $k+\lambda/\delta$; a chained
immediate-transfer process with negative-binomial rupture sizes and interval Laplace
transforms. &
The branching property is classical; the yield identity, and the failure of the naive free
sum, are the additions. \\
\addlinespace[3pt]

\anchor{Def 6.1, Eq.\ 6.11}\newline pp.\,153--155 &
An age-structured within-host system driven by \emph{closed-form} birth--death--catastrophe
kernels, with no free production kernel left to fit. &
Nelson, Gilchrist, Coombs, Hyman and Perelson (2004), who leave the age-dependent kernel
$P(a)$ free. \\
\addlinespace[3pt]

\anchor{Thm 6.3}\newline p.\,157 &
A projection of the exponential phase onto an effective pair
$(p_{\mathrm{eff}},d_{I,\mathrm{eff}})$, with closed endpoints --- so a fitted burst rate
is a property of the phase, not of the cell. &
Euler--Lotka renewal theory; Wallinga and Lipsitch on generation intervals. \\
\addlinespace[3pt]

\anchor{\S6.5.6}\newline p.\,161 &
A concrete three-to-two identifiability obstruction: $(\lambda,\mu,\delta)$ is not
recoverable from population data alone. &
Miao, Xia, Perelson and Wu (2011) on structural non-identifiability. \\
\addlinespace[3pt]

\anchor{Prop 7.2}\newline p.\,179 &
A Feynman--Kac representation for the joint two-type occupation hazard, with a mesh-limit
argument for the pathwise exponential weight. &
Standard first-step decomposition for one type; the joint two-type object is the increment. \\
\addlinespace[3pt]

\anchor{Prop 7.9}\newline p.\,198 &
At $\lambda_2=0$ the solution is built on \emph{ordinary} Bessel functions
$J_\rho,Y_\rho$, a different boundary of the parameter space. &
Antal and Krapivsky's bi-critical case uses \emph{modified} Bessel functions, from a
different equation. \\
\addlinespace[3pt]

\anchor{Def 8.1, Prop 8.3}\newline pp.\,210--212 &
A replicator--suppressor process in which suppressors are consumed one-to-one and never
replenished. The escape wall $i>j$ is then rate-\emph{independent}: a counting argument, no
rates enter. &
Pilyugin and Antia (2000) limit clearance by \emph{handling time}: an engaged effector is
occupied, then returns to the free pool. Their threshold therefore depends on the rates. \\
\addlinespace[3pt]

\anchor{Prop 8.7, Cor 8.8}\newline pp.\,218--219 &
Under an integer removal budget $\vartheta$, the surviving fraction is
$\mathcal{R}(\vartheta)=V_\infty a^{-\vartheta}$ with threshold
$\vartheta^\ast=\log V_\infty/\log a$. &
Geometric stop-loss is actuarial folklore; not found for a catastrophe process under an
integer immune budget. \\
\addlinespace[3pt]

\anchor{Prop 8.12}\newline p.\,220 &
Under threshold removal the flooding criterion \emph{reverses}: bursting now wins exactly
when $L<1$. The same dimensionless group orders the comparison in both directions. &
Yuan and Allen (2011) find the ranking of budding and bursting to depend on whether the immune
response is activated, by a different mechanism. \\
\addlinespace[3pt]

\anchor{Cor 9.3}\newline p.\,233 &
A two-layer clock: speciation intensity tracks unused \emph{individual} capacity through a
separate logistic abundance layer, not the species count. The limit law is geometric, with
a variance not previously recorded. &
Kendall (1948) time-dependent Yule; diversity-dependent models act on species number. \\
\addlinespace[3pt]

\anchor{Eq.\ 9.21}\newline p.\,240 &
A closed form for the moving unstable critical point of the Pimentel resistance chain: a
reversal is that point crossing the leading population. &
Pimentel (1965) supplies the chain; no antecedent was identified for the critical point. \\
\end{xltabular}\par}

% =====================================================================
\oisec{3. Inherited, and cited as such}

{\footnotesize\noindent
Naming these is what makes the two sections above legible as increments.
\textbf{Karlin and Tavar\'e (1982)} supply the killed birth--death process itself: the
linear-fractional generating function, the closed forms $I(t)$, $D(t)$, $J(t)$, $K(t)$, the
roots $a,b$, the geometric occupancy and the geometric killing position. \textbf{Brockwell,
Gani and Resnick (1982)} is the more general neighbour, admitting immigration and random
catastrophe sizes. \textbf{Kolmogorov (1938)} and \textbf{Yaglom (1947)} give the constant
$A(p)$, the critical decay $S_n\sim2/n$ and the conditioned limit law; \textbf{Pakes (2026)}
is the current word on whether that constant admits an explicit form.
\textbf{Kendall (1948)} gives the linear birth--death generating function reused throughout,
and the time-dependent Yule process of Chapter~9. \textbf{Antal and Krapivsky (2011)} give
the triangular two-type process and its Gauss reduction. \textbf{Perelson (1996)}, \textbf{Nowak
and Bangham (1996)} and \textbf{McLean (1993)} give constant-rate viral replication;
\textbf{Pearson, Krapivsky and Perelson (2011)} give the invariance of $R_0$ under matched
lifetime yield and the establishment comparison; \textbf{Komarova (2007)} gives flooding as a
biological idea, and \textbf{Yuan and Allen (2011)} a published ranking switch under an immune
term. \textbf{Nelson et al.\ (2004)} give the age-structured system; \textbf{Miao et al.\
(2011)} non-identifiability; \textbf{Pilyugin and Antia (2000)} finite nonspecific killers.
\textbf{Nee, May and Harvey (1994)} give the push of the past, against a large
diversity-dependent literature. \textbf{Ke, Chen and Yang (2013)} give the temperature- and
pH-dependent account of the \yp{} switch.\par}

% =====================================================================
\oisec{4. Scope of the claims}

{\footnotesize\noindent
``Not found'' is not a proof of absence. Where the evidence was thin the claim was weakened
rather than strengthened, and every statement in Section~1 has a published ancestor supplying
the method --- so each is written as an original identity obtained by specialising a known
technique. Three things are deliberately \emph{not} claimed: the killed birth--death process
and its closed forms, which are Karlin and Tavar\'e's; the constant $A(p)$, which is
Kolmogorov's; and the comparison of bursting with budding as a question, which is Pearson's
and Komarova's. What is claimed is the calculus built on the first, the Koenigs dictionary
for the second, and the criterion $L$ that settles the third in both directions.\par}

% =====================================================================
\oisec{5. Quick reference: where the contributions sit}

\vspace{0.15\baselineskip}
{\footnotesize
\begin{xltabular}{\linewidth}{@{}p{7mm}p{15mm}Z@{}}
\toprule
\oihead{Ch} & \oihead{Pages} & \oihead{What it contributes} \\
\midrule
\endfirsthead
\toprule
\oihead{Ch} & \oihead{Pages} & \oihead{What it contributes} \\
\midrule
\endhead
\bottomrule
\endlastfoot
1  & 7--21    & Framing, biology and motivation. No result is claimed here. \\
2  & 22--57   & Toolkit: Markov chains, branching processes, the birth--death--catastrophe process. The increment is gated release with a state-dependent opening intensity (p.\,35). \\
3  & 58--87   & Kolmogorov's constant $A(p)$, recast as the Koenigs linearising function of the logistic map and shown hypertranscendental (\anchor{Thm 3.5}, p.\,82). \\
4  & 88--111  & The intracellular calculus: the released count $W_t$ as a second coordinate, and every load moment as a polynomial in $I$ (\anchor{Thms 4.13--4.14}, \anchor{Prop 4.17}). \\
5  & 112--147 & The complete law of one bursting cell --- burst size, burst timing, several founders, chained transfer (\anchor{Thm 5.10}, \anchor{Props 5.12--5.13}). \\
6  & 148--175 & \textbf{Two of the four new statements.} The renewal embedding, and the criterion $L$ ordering bursting against budding (\anchor{Thm 6.8}, \anchor{Prop 6.9}). \\
7  & 176--204 & \textbf{Two of the four new statements.} Exact finite-time survival under two-type killing, and the killed bivariate generating function (\anchor{Thm 7.4}, \anchor{Eq.\ 7.50}). \\
8  & 205--227 & Two finite-budget immune caricatures, and the reversal of the flooding criterion under threshold removal (\anchor{Prop 8.12}). \\
9  & 228--256 & A two-layer capacity clock on a Yule process (\anchor{Cor 9.3}) and a moving critical point in the Pimentel chain (\anchor{Eq.\ 9.21}). \\
10 & 257--272 & Synthesis, and a consistency check on the \yp{} release strategy (\anchor{Eq.\ 10.9}, p.\,263) offered explicitly as a check rather than a proof. \\
\end{xltabular}\par}
